\documentclass[acmsmall,screen,review,anonymous]{acmart}

\settopmatter{}

\clubpenalty = 10000
\widowpenalty = 10000
\displaywidowpenalty = 10000

%% Bibliography style
\bibliographystyle{ACM-Reference-Format}
%% Citation style
%% Note: author/year citations are required for papers published as an
%% issue of PACMPL.
\citestyle{acmauthoryear}   %% For author/year citations
%% \citestyle{acmnumeric}
\usepackage{mathpartir}
\usepackage{tikz-cd}
\usepackage{lipsum}
\usepackage{enumitem}
\usepackage{wrapfig}
\usepackage{fancyvrb}
\usepackage{xspace}
\usepackage{thmtools} % needed for sane cleveref naming of theorems
\usepackage{mathtools}
\usepackage{hyperref}
\usepackage[capitalise]{cleveref}
\usepackage[dvipsnames]{xcolor}
\usepackage{minted}
% \PassOptionsToPackage{prologue,dvipsnames}{xcolor}

% \usepackage{showframe}

\usepackage[LGR,T1]{fontenc}
\DeclareMathAlphabet{\mathgtt}{LGR}{cmtt}{m}{n}

% \renewcommand{\familydefault}{\ttdefault}
% \usepackage{mathastext}
% \renewcommand{\familydefault}{\rmdefault}

\colorlet{keywordcolor}{blue}
\colorlet{nonlintycolor}{RubineRed}
\colorlet{lintycolor}{PineGreen}
\colorlet{nonlinconstructorcolor}{nonlintycolor}
\colorlet{linconstructorcolor}{lintycolor}
\colorlet{alphabetcolor}{black}
\colorlet{funarrowcolor}{lintycolor}
\definecolor{backcolour}{rgb}{0.95,0.95,0.92}

\newcommand{\uparrowcode}{\color{nonlinconstructorcolor} \uparrow}
\newcommand{\tocode}{\color{nonlinconstructorcolor} \to}
\newcommand{\tocodelin}{\color{linconstructorcolor} \to}

\usepackage{listings}
\lstdefinelanguage{Agda}{
alsoletter={0123\&'},
keywords={data, where, module, import, open, public,
record, field, let, in, if, then, else, case, of, with,
do, postulate, primitive, mutual, abstract, private,
forall, exists, cong, set, prop, sort, Level, Data, Type,
Renamer},
morekeywords=[2]{L, Bool, Nat, Fin},
morekeywords=[3]{String, Dyck, Trace, Exp, Atom, O, D, C, A, NotStartsWithLP, NotStartsWithRP},
morekeywords=[4]{bal, .nil, .cons, cons, nil, inl, inr, stop, 1to1, 1to2, 0to2, 0to1,
  done, add, num, parens, left, unexpected, lookAheadRP, lookAheadNot, closeGood, closeBad,
  doneGood, doneBad},
morekeywords=[5]{.states, states, .isAcc, isAcc, transitions, src, dst, .transitions, .dst,
  .src, .label, label, 0 , 1 , 2, 3, true, false},
morekeywords=[6]{NUM},
morekeywords=[7]{\&, '},
sensitive=true,
comment=[l]{--},
morecomment=[s]{\{-}{-\}},
morestring=[b]",
mathescape=true,
moredelim={[is][]{\#}{\#}},
escapeinside={(*@}{@*)},
escapebegin=\color{funarrowcolor},
escapeend={}
}
\lstset{
language=Agda,
basicstyle=\ttfamily\footnotesize,
keywordstyle=\color{keywordcolor},
keywordstyle=[2]\color{nonlintycolor},
keywordstyle=[3]\color{lintycolor},
keywordstyle=[4]\color{linconstructorcolor},
keywordstyle=[5]\color{nonlinconstructorcolor},
keywordstyle=[6]\color{alphabetcolor},
keywordstyle=[7]\color{funarrowcolor},
identifierstyle=\color{black},
commentstyle=\color{gray}\textit,
stringstyle=\color{orange},
numbers=none,
numberstyle=\tiny\color{gray},
stepnumber=1,
numbersep=5pt,
showstringspaces=false,
tabsize=4,
captionpos=b,
breaklines=true,
breakatwhitespace=false,
rulecolor=\color{black},
texcl=true,
backgroundcolor=\color{backcolour},
frame=single,
xleftmargin = 1.5mm ,
xrightmargin = 1.5mm
}

\usepackage{placeins} % for \FloatBarrier

\lstnewenvironment{floatlisting}{
  \FloatBarrier
  \lstset{
    aboveskip=2mm
  }
}{}

\usepackage{stmaryrd}
\usepackage{scalerel}
\usepackage{tikz}

\usepackage{pdfpages}

\newif\ifdraft
\drafttrue
\newcommand{\steven}[1]{\ifdraft{\color{orange}[{\bf Steven says}: #1]}\fi}
\renewcommand{\max}[1]{\ifdraft{\color{blue}[{\bf Max says}: #1]}\fi}
\newcommand{\johnson}[1]{\ifdraft{\color{red}[{\bf Johnson says}: #1]}\fi}
\newcommand{\pipe}{\,|\,}

\newcommand{\states}{\mathtt{states}}
\newcommand{\labelt}{label}
\newcommand{\transitions}{\texttt{transitions}}
\newcommand{\epstransitions}{\epsilon\texttt{transitions}}
\newcommand{\isAcc}{\mathtt{isAcc}}
\newcommand{\init}{\mathtt{init}}
\newcommand{\src}{\mathtt{src}}
\newcommand{\dst}{\mathtt{dst}}
\newcommand{\pparse}{\mathtt{parse}}
\newcommand{\print}{\mathtt{print}}
\newcommand{\epssrc}{\epsilon\mathtt{src}}
\newcommand{\epsdst}{\epsilon\mathtt{dst}}
\newcommand{\balanced}{\mathtt{balanced}}
\newcommand{\Dyck}{\mathtt{Dyck}}
\newcommand{\N}{\texttt{N}}
\newcommand{\Cl}{\mathsf{Cl}}

\newcommand{\iso}[2]{#1 \cong #2}
\newcommand{\id}{\mathsf{id}}
\newcommand{\isoover}[3]{#1 \cong_{#2} #3}
\newcommand{\define}[2]{#1 \coloneqq #2}
\newcommand{\prn}[1]{\left( #1 \right)}
\newcommand{\cat}[1]{\mathcal{#1}}
\newcommand{\inob}[2]{#1 : \mathcal{#2}}
\newcommand{\optext}{\mathsf{op}}
\newcommand{\opsuperscript}[1]{#1^{\optext}}
\newcommand{\op}[1]{\opsuperscript{\cat{#1}}}
\newcommand{\seqstar}[2]{#1 \mathbin{\star} #2}
\newcommand{\seqcirc}[2]{#2 \mathbin{\circ} #1}
\newcommand{\seq}[2]{\seqstar{#1}{#2}}
\newcommand{\cathom}[3]{\cat{#1}\prn{#2,#3}}
\newcommand{\cathomraw}[3]{{#1}\prn{#2,#3}}
\newcommand{\incathom}[4]{#1 : \cathom{#2}{#3}{#4}}
\newcommand{\incathomraw}[4]{#1 : \cathomraw{#2}{#3}{#4}}
\newcommand{\dispsuperscript}[1]{#1^{\mathbf{D}}}
\newcommand{\vertsuperscript}[1]{#1^{\mathbf{V}}}
\newcommand{\mkdisp}[1]{\dispsuperscript{#1}}
\newcommand{\disp}[1]{\dispsuperscript{\mathcal{#1}}}
\newcommand{\disppsh}[1]{\dispsuperscript{#1}}
\newcommand{\vertpsh}[1]{\vertsuperscript{#1}}
\newcommand{\catofelts}[1]{\mathsf{Elt}\prn{#1}}
\newcommand{\dispob}[2]{\disp{#1}_{#2}}
\newcommand{\dispobraw}[2]{{#1}_{#2}}
\newcommand{\indispob}[3]{#1 : \dispob{#2}{#3}}
\newcommand{\disphom}[4]{\disp{#1}_{#2}\prn{#3,#4}}
\newcommand{\disphomraw}[4]{{#1}_{#2}\prn{#3,#4}}
\newcommand{\indisphom}[5]{#1 : \disphom{#2}{#3}{#4}{#5}}
\newcommand{\dispopsuperscript}[1]{{\prn{#1}}^{\dispsuperscript{\optext}}}
\newcommand{\dispop}[1]{\dispopsuperscript{\disp{#1}}}
\newcommand{\reindex}[2]{\mathrm{reindex}_{#1}\prn{#2}}
\newcommand{\totalcat}[1]{\textstyle\int #1}
\newcommand{\vertV}[1]{#1^{\mathsf{V}}}

\newcommand{\Set}{\mathsf{Set}}
\newcommand{\SetD}{\dispsuperscript{\Set}}
\newcommand{\Psh}{\mathsf{Psh}}
\newcommand{\PshD}{\dispsuperscript{\Psh}}
\newcommand{\hSet}{\mathsf{hSet}}
\newcommand{\Type}{\mathsf{Type}}

\newcommand{\objtxt}{\mathsf{obj}}
\newcommand{\elttxt}{\mathsf{elt}}
\newcommand{\obj}[1]{#1_{\objtxt}}
\newcommand{\elt}[1]{#1_{\elttxt}}
\newcommand{\pshelt}[2]{#1~#2}
\newcommand{\dispobj}[1]{#1_{\mkdisp{\objtxt}}}
\newcommand{\dispelt}[1]{#1_{\mkdisp{\elttxt}}}
\newcommand{\disppshelt}[3]{\disppsh{#1}\prn{#2, #3}}

\newcommand{\PiTy}[3]{\textstyle\prod_{#1 : #2} #3}
\newcommand{\SigTy}[3]{\textstyle\sum_{#1 : #2} #3}
\newcommand{\PiTyLimit}[3]{\textstyle\prod\limits_{#1 : #2} #3}
\newcommand{\SigTyLimit}[3]{\textstyle\sum\limits_{#1 : #2} #3}

\newcommand{\Mon}[1]{\mathrm{Mon}[#1]}
\newcommand{\CCCtxt}{\mathrm{CCC}}
\newcommand{\CCC}[1]{\CCCtxt[#1]}
\newcommand{\BiCCC}[1]{\mathrm{BiCCC}[#1]}
\newcommand{\CCtxt}{\mathrm{CC}}
\newcommand{\CC}[1]{\CCtxt[#1]}
\newcommand{\lang}{\mathcal{L}}
\newcommand{\langsig}[1]{\lang [#1]}

\newcommand{\dom}{\mathsf{dom}}
\newcommand{\cod}{\mathsf{cod}}
\newcommand{\rec}{\mathsf{rec}}
\newcommand{\recinterp}[1]{\rec_{#1}}
\newcommand{\recinterpCC}[1]{\rec^{\CCtxt}_{#1}}
\newcommand{\recinterpCCC}[1]{\rec^{\CCCtxt}_{#1}}
\newcommand{\elim}{\mathsf{elim}}
\newcommand{\eliminterp}[1]{\elim_{#1}}
\newcommand{\bool}{\mathsf{Bool}}
\newcommand{\boolt}{\mathsf{true}}
\newcommand{\boolf}{\mathsf{false}}
\newcommand{\List}{\mathsf{List}}
\newcommand{\nil}{\mathsf{nil}}
\newcommand{\cons}{\mathsf{cons}}
\newcommand{\ListOf}[1]{\List~#1}
\newcommand{\reclist}{\rec_{\List}}
\newcommand{\elimlist}{\elim_{\List}}

\usepackage{CJKutf8}
\newcommand{\yochar}{\text{\begin{CJK}{UTF8}{min}よ\end{CJK}}}
\newcommand{\yo}[1]{\yochar\,#1}
\newcommand{\nervetxt}{N}
\newcommand{\nerve}[1]{\nervetxt_{#1}}
\newcommand{\dispyochar}{\dispsuperscript{\yochar}}
\newcommand{\dispyo}[1]{\dispyochar\,#1}
\newcommand{\recyo}[1]{\rec_{\yochar}\prn{#1}}
\newcommand{\recdispyo}[1]{\rec_{\dispyochar}\prn{#1}}
\newcommand{\IsoFiber}{\mathrm{IsoFiber}}
\newcommand{\Iso}{\mathrm{Iso}}

\newcommand{\cartlift}[2]{#1^{*} #2}
\newcommand{\cartliftbare}[1]{#1^{*}}
\newcommand{\vertprod}[2]{#1 \times^{\mathbf{V}} #2}
\newcommand{\dispprod}[2]{#1 \times^{\mathbf{D}} #2}
\newcommand{\vertexp}[2]{#1 \Rightarrow^{\mathbf{V}} #2}
\newcommand{\dispexp}[2]{#1 \Rightarrow^{\mathbf{D}} #2}
\newcommand{\uq}[2]{\forall_{#1}\,#2}
\newcommand{\lamforall}[1]{\lambda^{\forall}_{#1}}
\newcommand{\appforall}[1]{\mathsf{app}^{\forall}_{#1}}
\newcommand{\wk}[1]{\mathsf{wk}_{#1}}
\newcommand{\app}{\mathsf{app}}

\newcommand{\pext}[1]{\mathrm{ext}\prn{#1}}
\newcommand{\coy}[1]{\diamond #1}

\newcommand{\slicecat}[2]{#1 / #2}

\newcommand{\nat}{\mathsf{Nat}}

\newcommand{\isTy}[1]{#1\;\mathsf{type}}
\newcommand{\sem}[1]{\llbracket #1 \rrbracket}
\newcommand{\interp}{\leadsto}
\newcommand{\interpIn}[2]{#1 \leadsto #2}
\newcommand{\Agda}{\raisebox{-0.2ex}{\includegraphics[height=1.8ex]{logo-agda.pdf}}\xspace}

\newcommand{\extref}[3]{#1~#2 in~\cite{#3}}

% These declarations are to get cleveref to play nice with
% different theorem-like environments. Otherwise \cref{lem1.1}
% will display 'Theorem 1.1'

\declaretheoremstyle[headfont=\normalfont\itshape]{defstyle}
\declaretheoremstyle[headfont=\normalfont\scshape]{thmstyle}

\declaretheorem[name=Definition, style=defstyle, numberwithin=section]{definition}
\declaretheorem[style=thmstyle, name=Theorem, numberlike=definition]{theorem}
\declaretheorem[style=thmstyle, name=Axiom, numberlike=definition]{axiom}
\declaretheorem[style=thmstyle, name=Corollary, numberlike=definition]{corollary}
\declaretheorem[style=thmstyle, name=Construction, numberlike=definition]{construction}

\crefname{definition}{Definition}{Definitions}
\crefname{theorem}{Theorem}{Theorems}
\crefname{axiom}{Axiom}{Axioms}
\crefname{corollary}{Corollary}{Corollaries}
\crefname{construction}{Construction}{Constructions}

\begin{document}

\pagebreak

\title{A Call-by-push-value Scheme}

\author{Max S. New}
\orcid{0000-0001-8141-195X}
\affiliation{%
  \institution{University of Michigan}
  \city{Ann Arbor}
  \country{USA}
}
\email{maxsnew@umich.edu}

\begin{abstract}
Scheme is known for its minimalism: what would be core language features in other languages are instead implemented as macros over a small kernel. In its most Spartan formulation all compound data arises from the humble cons. Inspired by Levy's Call-by-push-value, we propose a variant of Scheme that applies that same minimalism to multi-argument functions. That is, rather than having both lambda and lambda dot args, we have a single primitive for checking whether there are any more arguments on the stack. In this view, function application is interpreted as a "cons on the stack". We provide a proof-of-concept implementation called Fiddle implemented in Racket, and we discuss other possible applications of this dual viewpoint.
\end{abstract}

%% 2012 ACM Computing Classification System (CSS) concepts
%% Generate at 'http://dl.acm.org/ccs/ccs.cfm'.

\maketitle

\section{Making Scheme more Minimal}

Scheme is known for minimalism. A small kernel language includes few
primitives out of which more complex language features are built up
using macros \cite{growing-a-language}.

In an extremely minimalistic version of Scheme, the only types of data
would be atoms, cons cells and closures. Inductive datatypes like
lists are not primitives, but are emergent from the combination of
cons cells and a single atom \mintinline{racket}{nil}.
%
Rather than being a core primitive, pattern-matching syntax can be
built up using macros which call out to the basic constructs:
\mintinline{racket}{nil?} \mintinline{racket}{cons?} \mintinline{racket}{car} and \mintinline{racket}{cdr} as well as
\mintinline{racket}{if} expressions.

But even minimalistic Schemes typically do include a ``compound''
construct: closures can take any number of arguments, and these can be
accessed \emph{as a list}. For example to define a function that
returns its first argument, no matter how many additional arguments
are passed:
\begin{minted}{racket}
  (define first-arg (x . l) x)
\end{minted}

This is mirrored in the process of \emph{calling} functions:
\begin{minted}{racket}
  (f w x y z)
  (apply f l)
\end{minted}
Function application is a family of many-argument operations with the
universal way to perform function application being the apply function
which takes a list as an argument.

This presents a kind of asymmetry: data types in Scheme are built up
using only simple constructors, but functions, i.e.,
\emph{codata} types build in the use of a compound inductive structure.

\section{Fiddle: A CBPV Scheme}

We propose a language, \emph{Fiddle}, which brings the same
minimalistic philosophy of Scheme datatypes to its codatatypes. The
approach is based on Levy's call-by-push-value calculus (CBPV)
\cite{cbpv}, and in particular on a semantic model of dynamically
typed CBPV developed in the study of gradual typing
\cite{gradual-cbpv}.

The core idea of Fiddle's CBPV approach to functions is to recognize
that in Scheme and other call-by-value languages, a function call is really a complex construct. In the following call to \mintinline{racket}{f}:
\begin{minted}{racket}
  (+ 1 (f x y))
\end{minted}
Three distinct operations are performed:
\begin{enumerate}
\item A continuation is created from the context \mintinline{racket}{(+ 1 -)}
\item The arguments \mintinline{racket}{x y} are supplied
\item Control is transferred to \mintinline{racket}{f}
\end{enumerate}
In Fiddle, these are split into three separate notions, which we
introduce in the opposite order.

First, closures are a primitive datatype just as in Scheme, but the primitive form for creating a closure is not a lambda but the suspension form \mintinline{racket}{(~ e)} which creates a closure that can be invoked using the force form \mintinline{racket}{(! v)}. Forcing a suspended closure runs the suspended code, but the closure itself is a first-class value.

Secondly, arguments are supplied through a single argument form \mintinline{racket}{(e v)} which operationally has the effect of \emph{pushing} the value \mintinline{racket}{v} onto the current stack and then executing the expression \mintinline{racket}{e}. A computation \emph{pops} an argument off of the stack using the lambda form \mintinline{racket}{(lambda (x) e)}. So in Fiddle (and CBPV more generally), the lambda form is not itself a value but instead an imperative operation that pops the value from the stack -- the origin of the term call-by-push-value. A multi-argument call is then syntactic sugar for a sequence of pushes, so \mintinline{racket}{(! f x y)} desugars to \mintinline{racket}{(((! f) x) y)}.

Lastly, continuations in Fiddle are created through an explicit
sequencing operation resembling a monadic let-binding
\mintinline{racket}{(bind (x e1) e2)}. Operationally, this creates a
continuation that captures the current stack, and invokes
\mintinline{racket}{e1} with a now empty stack of arguments. In the
current implementation of fiddle, returning is an explicit operation
\mintinline{racket}{(ret v)} which invokes the current continuation
with the value \mintinline{racket}{v}.

Combined together the program is written in Fiddle as
\begin{minted}{racket}
  (bind (r (! f x y))
    (! + 1 r))
\end{minted}
The separation of the creation of a continuation from the invocation
of a procedure closely resembles Steele's classic explanation of how to implement procedure calls \cite{steele-goto}.

%% \begin{minted}{racket}
%%   (f w x y z)
%% \end{minted}
%% Into a sequence of one-argument applications, separate from a construct for actually transferring control:
%% \begin{minted}
%%   (((((! f) w) x) y) z)
%% \end{minted}
%% The one argument application form \mintinline{racket}{(e x)} has the effect of \emph{pushing} the \mintinline{racket}{x} onto the stack, and then executing \mintinline{racket}{e}. The jump form \mintinline{racket}{(! f)} invokes a closure, effectively a tail call where the arguments are passed as the implicit state of the stack.

%% On the callee side, we need two corresponding forms. First \mintinline{racket}{(~ e)} \emph{suspends} the computation \mintinline{racket}{e}
%% \begin{minted}{racket}
%%   (~ e)
%% \end{minted}

\section{Implementation}

We implemented Fiddle as an embedded language in Racket in 2020. We
used the \emph{turnstile} library \cite{turnstile} for defining typed
languages, with the very minimal type system of only two syntactic
sorts: values and computations. We have implemented a small library,
including a complex copattern matching macro to make defining variable
arity functions more ergonomic, as well as lazy datastructures and
abstract datatypes using Fiddle's codata types.

To implement the CBPV-style of stack, the Fiddle runtime has a single
mutable list which is mutated by pushing and popping. To facilitate
reuse of Racket's standard libraries we also implement wrapper
transformations that allow Racket functions to be called from Fiddle
and vice-versa, transferring values from the Fiddle stack to the
Racket stack and vice-versa.

We used Fiddle as a prototype for CBPV programming, implementing some
programming puzzles such as Advent of code, which led to implementing
lexers, some basic data structures and monad-like abstractions for
effects. This has since led to the development of Zydeco, a typed CBPV
language and formalized the monad abstraction used in Fiddle as
relative monads \cite{zydeco, relative-monads}.

\bibliography{refs.bib}

\end{document}
